\section{Comparación unificada de los tres métodos}

\begin{table}[H]
\centering
\small
\begin{tabular}{>{\raggedright\arraybackslash}p{2.2cm}>{\raggedright\arraybackslash}p{3.4cm}>{\raggedright\arraybackslash}p{3.4cm}>{\raggedright\arraybackslash}p{4.2cm}}
\toprule
\textbf{Método} & \textbf{Etapa principal} & \textbf{Representación central} & \textbf{Beneficio / costo principal} \\
\midrule
LATS & Búsqueda en tiempo de inferencia & Árbol de estados y acciones & Permite retroceder y explorar varios planes; costo de llamadas alto \\
SPRINT & Inferencia paralela posentrenamiento & DAG de dependencias de tareas & Reduce la ruta crítica; identificar dependencias y agregar resultados es complejo \\
SWiRL & Datos fuera de línea + RL & Trayectorias multipaso con herramientas & Aprende un procedimiento general; limitado por el teacher, el judge y la cobertura fuera de línea \\
\bottomrule
\end{tabular}
\end{table}

\subsection{Formas de combinación}

Los tres métodos no se excluyen entre sí. Un sistema completo puede usar SPRINT para generar el grafo de tareas paralelas, ejecutar una búsqueda LATS dentro de cada subtarea difícil y luego usar la policy entrenada con SWiRL para elegir herramientas y procesar la observation.

\subsection{Elección de ingeniería}

\begin{itemize}
    \item Si el valor de la tarea es alto, es simulable y requiere retroceder: priorizar la búsqueda en árbol;
    \item si las subtareas son independientes y la latencia de las herramientas es alta: priorizar la ejecución paralela;
    \item si hay muchos entornos verificables y se busca una mejora a largo plazo: construir synthetic trajectories para hacer RL;
    \item si las acciones son irreversibles o el entorno es costoso: limitar la exploración en línea y aumentar el ensayo fuera de línea y la aprobación humana.
\end{itemize}

\subsection{Resumen de la sección}

El problema de la planeación puede resolverse en tres niveles: búsqueda, programación y aprendizaje. La búsqueda en árbol mejora la calidad de las decisiones, la programación paralela reduce el tiempo y el entrenamiento a nivel de proceso consolida la experiencia en la policy. El objetivo común de los tres es reducir la acumulación de errores y el cómputo inútil en trayectorias largas.
